\documentclass[11pt]{article}

\usepackage[margin=1in]{geometry}
\usepackage[T1]{fontenc}
\usepackage[utf8]{inputenc}
\usepackage{lmodern}
\usepackage{enumitem}
\usepackage{hyperref}
\usepackage{listings}

\hypersetup{colorlinks=true,linkcolor=blue,urlcolor=blue}

\lstset{basicstyle=\ttfamily\small,breaklines=true,columns=fullflexible,keepspaces=true,frame=single}

\title{Reference Comparison Experiment}
\author{}
\date{}

\begin{document}
\maketitle

\section{Purpose}

The reference comparison is a supplementary experiment added to connect the project to the Tang (2021) Geometric A* reference study. It does not replace the main experiment and it does not modify the main experiment protocol. The main experiment remains the thirty-two-configuration comparison between ECBS with classical mapping and ECBS with cyclic mapping.

The reference comparison now uses three separate 50x50 Tang-inspired port maps stored in the project:

\begin{lstlisting}
dev/inputs/reference_port_maps/port_map_1.png
dev/inputs/reference_port_maps/port_map_2.png
dev/inputs/reference_port_maps/port_map_3.png
\end{lstlisting}

These reference images are normalized to two ordinary cell colors: pure black pixels are obstacles and pure white pixels are free spaces. The former light-red ``invisible obstacle'' pixels have been converted to ordinary white/free cells, and the reference workflow no longer has a special invisible-obstacle pathfinding or visualization rule. The internal binary matrix follows the manuscript convention, with $1$ for traversable/free cells and $0$ for blocked/obstacle cells. The older orientation-based variants are no longer generated. The program treats the three input files as the three map cases directly.

The input directory may still contain sibling images with manually colored path markings, but those images are kept only as archived inputs. The reference-comparison program does not read those markings and does not import any manually traced path.

\section{Experiment Selector}

The project has a top-level selector in \texttt{dev/main.py}. Exactly one experiment family should be active:

\begin{lstlisting}
SELECTED_EXPERIMENT = "ref_comparison"
# SELECTED_EXPERIMENT = "main_experiment"
\end{lstlisting}

The reference-comparison configuration is separated from the main experiment configuration:

\begin{lstlisting}
dev/master_config_ref_comparison.py
\end{lstlisting}

This keeps the main \texttt{master\_config.py} untouched.

The reference comparison now uses the same cyclic-map post-processing variant as the main experiment:

\begin{lstlisting}
REMOVE_EXTRA_TRANSITIONS = True
ADD_TRANSITIONS_BETWEEN_FREE_SPACES = False
\end{lstlisting}

Redundant bidirectional transitions are reduced before required connectivity is restored, and no final all-free-space transition expansion is applied. The reference experiment therefore no longer uses a separate cyclic-map post-processing variant.

\section{Reference Workflow}

The reference workflow intentionally mirrors the main experiment workflow. The raw solver execution phase is separated from graph/data generation and visualization generation.

\begin{lstlisting}
# to_generate = "raw_data"
to_generate = "graphs"
# to_generate = "visualization"
\end{lstlisting}

If \texttt{to\_generate = "raw\_data"}, the selected reference case is executed, numerical raw data is saved, and the designated successful trajectories are written separately under \texttt{outputs\_ref\_comparison/frame\_by\_frame}. The single-agent case stores the first successful timing repetition for classical and cyclic in each of the three maps. The multi-agent case stores one final successful trajectory for classical at its discovered classical capacity and one for cyclic at its discovered cyclic capacity. Capacity-search trials are not stored as frame-by-frame files.

If \texttt{to\_generate = "graphs"}, the numerical raw data is reused. If \texttt{to\_generate = "visualization"}, the solver is not executed and the numerical raw data is not read; the visualizer loads the saved frame-by-frame packages and generates the Pillow frames from them. The visualization phase logs the current map and mapping, frame-writing stages, and final frame counts.

\section{Selectable Reference Cases}

The reference-comparison selector is:

\begin{lstlisting}
# SELECTED_PORT_EXPERIMENT = "single_agent"
SELECTED_PORT_EXPERIMENT = "multi_agent"
\end{lstlisting}

Both selectable cases use the same three 50x50 maps. The single-agent case runs the three maps with one agent. The multi-agent case performs separate classical and cyclic capacity searches for each map. Older selector names such as \texttt{single\_agent\_x50} and \texttt{multi\_agent\_x50} are kept only as backward-compatible aliases.

The multi-agent search controls are configured in \texttt{dev/master\_config\_ref\_comparison.py}:

\begin{lstlisting}
REFERENCE_CAPACITY_ATTEMPTS_PER_AGENT_NUMBER = 5
REFERENCE_CAPACITY_SUCCESSFUL_RUNS_REQUIRED = 3
REFERENCE_CAPACITY_PASS_CRITERION = "solver_success"
\end{lstlisting}

For each reference map, the candidate interval is $1\ldots F$, where $F$ is the number of traversable cells in that normalized map. Every tested agent number uses exactly five run slots and passes when at least three runs solve within the runtime limit. Classical and cyclic are judged independently. The binary search continues until the remaining numerical interval is resolved; there is no post-success early-stop move limit.

\section{Single-Agent Reference Case}

The single-agent reference case compares:

\begin{lstlisting}
Traditional A* with Classical Mapping
With Cyclic Mapping
\end{lstlisting}

Both single-agent sides are computed by the project using traditional A*. The classical side runs traditional A* on the classical map, while the cyclic side runs traditional A* on the cyclic map. Both use the same lower-left start cell and upper-right target cell on each of the three input maps. Each map and mapping is timed using five identical A* executions, controlled by \texttt{SINGLE\_AGENT\_TIMING\_REPETITIONS = 5}. The stored runtime is the average of those five timing samples.

The single-agent case records:

\begin{itemize}[leftmargin=2em]
    \item averaged time computation halted,
    \item A* nodes expanded,
    \item total path length,
    \item total turns.
\end{itemize}

When graph/data outputs are generated, the single-agent case also creates a matplotlib-formatted comparison table in the same graphs folder. In this table, Running time uses the averaged time computation halted value, Number of nodes uses the A* expanded-node count, Number of turns uses the total turns value, and Total distance uses the path length value. The first column is Map number, with sections for maps 1, 2, and 3. The table uses Traditional A* with Classical Mapping, With Cyclic Mapping, and Percent Reduction/Gain as the main comparison headers. The Path parameters header includes the italicized subtitle note \textit{(averages are across all agents)}. The percent column is signed: negative values mean the cyclic value is lower than the classical value, while positive values mean the cyclic value is higher. The raw run record also stores the individual timing samples used to compute the average.

The individual cyclic-faster filter is disabled for the single-agent case because cyclic mapping is not expected to help single-agent pathfinding. This case is included to show that cyclic mapping is not proposed as a general single-agent shortcut.

\section{Multi-Agent Reference Case}

The multi-agent reference case internally runs ECBS for both mappings. In the generated tables and graph legends, the two presented comparison labels are:

\begin{lstlisting}
Traditional A* with Classical Mapping
With Cyclic Mapping
\end{lstlisting}

For each port map, the multi-agent case searches classical capacity and cyclic capacity independently. A tested agent number passes for a mapping only when that mapping produces at least three successful solver runs out of five within the time limit. Runtime or conflict superiority over the other mapping is not part of the capacity definition, and the setup generation for one mapping does not require the other mapping to solve.

After both searches finish, classical mapping is evaluated at its discovered classical capacity and cyclic mapping is evaluated at its discovered cyclic capacity. All agents share the upper-right target cell and disappear after reaching it. Agents are released through two spawn cells: the upper neighbor and the right neighbor of the lower-left start cell. The solver supports explicit per-agent \texttt{spawn\_time}. Before an agent's spawn time, it does not occupy the map. At its spawn time, it appears at its assigned spawn cell.

The final comparison for each map and mapping is timed using three identical ECBS executions, controlled by \texttt{MULTI\_AGENT\_TIMING\_REPETITIONS = 3}. The stored runtime is the average of those three timing samples. The path-related values come from the representative deterministic solution.

The multi-agent case records:

\begin{itemize}[leftmargin=2em]
    \item averaged time computation halted,
    \item number of conflicts detected,
    \item total path length,
    \item total turns.
\end{itemize}

When graph/data outputs are generated, the multi-agent table uses Map number as the first column. The reported path parameters are Running time, Average number of conflicts, Average number of turns, and Average total distance. The table uses Traditional A* with Classical Mapping, With Cyclic Mapping, and Percent Reduction/Gain as the main comparison headers. The Path parameters header includes the italicized subtitle note \textit{(averages are across all agents)}. Per-agent quantities use the corresponding mapping's own discovered capacity as the divisor. The percent column is signed: negative values mean the cyclic value is lower than the classical value, while positive values mean the cyclic value is higher. Cyclic-faster filtering is not used in capacity search.

For both single-agent and multi-agent reference cases, graph outputs are now metric-level summaries across the three maps rather than one graph per map. Each graph uses Map number on the x-axis and contains four points: Map 1, Map 2, Map 3, and Average. The Average point is computed from the three map-level values for the same metric and mapping. Only Map 1, Map 2, and Map 3 are connected by the horizontal series line; the Average point remains a standalone marker so it is not visually treated as a fourth input map. The single-agent graphs cover Running time, Number of nodes, Number of turns, and Total distance. The multi-agent graphs cover Running time, Average number of conflicts, Average number of turns, and Average total distance. The graph legends use Traditional A* with Classical Mapping and With Cyclic Mapping.

\section{Heuristic and Solver Settings}

For both single-agent A* and multi-agent ECBS reference cases, low-level guidance uses Manhattan distance, matching the manuscript. Graph-aware true static shortest-path distance is disabled. The crowd-spreading/cohesion priority is also removed from the reference comparison.

Shared reference settings include:

\begin{lstlisting}
SHARED_TIME_LIMIT_SECONDS = 60.0
SHARED_ECBS_SUBOPTIMALITY = 3.0
SHARED_TRUE_STATIC_SHORTEST_PATH_DISTANCE = False
SHARED_TIGHT_TIME_HORIZON = False
\end{lstlisting}

\section{Output Layout}

Reference-comparison outputs are separate from the main experiment outputs. The root is:

\begin{lstlisting}
dev/outputs_ref_comparison/
\end{lstlisting}

Each case has its own independent folder:

\begin{lstlisting}
dev/outputs_ref_comparison/single_agent/
dev/outputs_ref_comparison/multi_agent/
\end{lstlisting}

Inside each case folder, the structure follows the same general idea as the main experiment:

\begin{lstlisting}
aggregates/
graphs/
logs/
metadata/
records/
visualizations/
\end{lstlisting}

Persisted raw reference data is stored under:

\begin{lstlisting}
dev/outputs_ref_comparison/raw_mapf_files/<case_id>/
\end{lstlisting}

For the multi-agent case, the raw payload and graph-record exports also include \texttt{capacity\_searches.json}. This file preserves every tested agent number, the five mapping-specific run slots, the success count, and the resulting pass or failure decision for both classical and cyclic searches. The per-map traversable-cell counts and the separately discovered classical and cyclic capacities are also written to \texttt{stop\_summary.json}.

\section{Important Interpretation Notes}

The reference comparison is supplementary. It should not be merged with the main thirty-two-configuration results. It is intended as a focused reference check using three 50x50 port maps, fixed lower-left to upper-right movement, and repeated timing samples to reduce fluke timing effects.

The main experiment remains the primary evidence for cyclic mapping's overall behavior.

\end{document}
